\chapter{L'autre côté de la palissade}
\label{chap:lautre-cote-de-la-palissade}

Kamatsu entendit les cris avant de comprendre ce qui se passait.

Il marchait entre deux chariots lorsque des voix s'élevèrent à l'avant de la
caravane. Un cheval hennit brusquement, bientôt suivi par un autre. Des hommes
coururent vers l'origine du tumulte tandis que d'autres criaient aux enfants
de rejoindre leurs familles.

Kamatsu resta quelques secondes à regarder avant que quelqu'un ne l'attrape par
le bras.

— Kamatsu !

Sa mère le ramena contre elle et l'entraîna vers leur chariot. Son père les
rejoignit presque aussitôt. Il tenait une arme que Kamatsu lui avait rarement
vue entre les mains.

— Reste avec ta mère.

— Qu'est-ce qui se passe ?

Son père ne répondit pas.

Autour d'eux, la caravane qu'il avait toujours connue semblait soudain ne plus
fonctionner comme elle aurait dû. Les chevaux ne suivaient plus la route. Les
chariots s'arrêtaient de travers ou se heurtaient presque. Les voix familières
n'échangeaient plus des plaisanteries ou des indications sur le voyage : elles
criaient des ordres, des avertissements et parfois des noms.

Kamatsu aperçut alors les Orcs.

Il en avait déjà entendu parler autour des feux, mais les histoires n'étaient
pas toujours d'accord entre elles. Certains voyageurs les présentaient comme
des guerriers féroces, d'autres comme des brutes que l'on pouvait entendre
arriver de loin. Kamatsu en avait imaginé de toutes sortes.

Ceux qui attaquaient la caravane n'avaient rien d'imaginaire.

Sa mère le poussa derrière une roue et s'accroupit avec lui. Kamatsu voulut
regarder par-dessus son épaule, mais elle posa une main derrière sa tête et le
ramena contre elle.

— Reste là.

Cette fois, il obéit.

Il entendit davantage qu'il ne vit : le choc du métal, le bois qui craquait,
les chevaux affolés, des pas qui couraient tout autour d'eux. À plusieurs
reprises, Kamatsu demanda où était son père, mais sa mère répondit seulement
qu'il allait revenir.

Et il revint.

Pas seul.

Deux Orcs l'accompagnaient. Il n'avait plus son arme. Kamatsu sentit les bras
de sa mère se resserrer autour de lui.

Après cela, tout alla très vite. On les fit sortir de derrière le chariot.
Kamatsu chercha les visages qu'il connaissait parmi ceux qui se trouvaient
encore sur la route. Il en reconnut plusieurs, regroupés sous la surveillance
des Orcs. D'autres manquaient.

Il voulut demander où ils étaient.

En voyant le visage de sa mère, il ne posa pas la question.

\scenebreak

Ils quittèrent la route sous bonne garde.

Kamatsu marcha entre ses parents, tenant la main de sa mère. Son père restait
de l'autre côté, assez près pour qu'il puisse le toucher lorsqu'il en
ressentait le besoin. Aucun des deux ne parlait beaucoup.

Kamatsu, lui, avait des questions.

— Où est-ce qu'ils nous emmènent ?

— Je ne sais pas.

Quelques minutes plus tard, il recommença.

— Ils habitent dans la forêt ?

— Peut-être.

Puis, après avoir encore marché quelque temps :

— On va revenir aux chariots ?

Cette fois, sa mère mit un peu plus de temps à répondre.

— Reste près de nous pour l'instant.

Kamatsu acquiesça et continua de marcher.

La forêt aurait pu l'occuper pendant des heures en d'autres circonstances. Il
y avait des arbres aux formes étranges, des racines qui sortaient du sol et
des trouées entre les branches dans lesquelles le soleil dessinait des taches
de lumière. À plusieurs reprises, son regard s'y attarda, mais il revenait
toujours aux Orcs qui les entouraient.

Il observa leurs armes, leurs vêtements, leur manière de se déplacer entre les
arbres. Certains parlaient entre eux dans une langue qu'il ne comprenait pas.
Kamatsu essaya d'en retenir quelques sons avant de renoncer et se rapprocha un
peu de son père.

Le voyage lui parut long.

Lorsque la forêt finit par s'ouvrir, Kamatsu découvrit une palissade dressée
sur une hauteur. Derrière elle apparaissaient des toits, des fumées et des
silhouettes qui se déplaçaient entre les constructions.

— C'est là qu'ils habitent ?

Son père regarda le village.

— Oui.

Kamatsu continua de l'observer tandis qu'on les conduisait vers l'entrée. Il
avait vu des villages entourés de murs au cours de ses voyages, mais celui-ci
ne ressemblait à aucun de ceux qu'il connaissait. Tout y paraissait plus
massif, plus rude. Les regards se tournaient vers eux à mesure qu'ils
avançaient.

Puis Kamatsu aperçut le trou.

Il ralentit malgré lui.

Au milieu du village, le sol s'ouvrait sur une excavation immense, assez large
pour que plusieurs maisons aient pu tenir à l'intérieur. Une solide barrière
de bois en suivait le contour, séparant les habitants du vide. De l'autre
côté, les parois plongeaient jusqu'à un vaste fond de terre battue, plusieurs
mètres plus bas.

Pendant quelques secondes, Kamatsu oublia presque les Orcs autour de lui. Il
n'avait jamais rien vu de semblable.

— Maman, regarde.

Sa mère suivit son regard.

— Qu'est-ce que c'est ?

Elle ne répondit pas tout de suite.

Kamatsu, lui, cherchait déjà une explication. Peut-être avait-on creusé pour
trouver quelque chose sous le village. Un trésor, songea-t-il aussitôt. Ou
peut-être existait-il un passage quelque part au fond, menant sous la colline.

Puis il remarqua quelque chose d'étrange. La barrière n'était pas seulement là
pour empêcher de tomber. Par endroits, le sol avait été piétiné par les
nombreux pieds qui s'étaient tenus devant elle.

Des gens venaient souvent regarder ce qui se passait là-dessous.

Kamatsu se rapprocha légèrement de la barrière pour essayer de mieux voir, mais
la main de sa mère se resserra autour de la sienne.

— Viens, Kamatsu.

Elle l'entraîna avec elle avant qu'il puisse s'approcher davantage. Kamatsu
tourna encore la tête en marchant, essayant d'imaginer ce qui pouvait bien
attirer autant de monde autour d'un trou pareil.

Pour la première fois depuis leur entrée dans le village, il ne posa pas la
question.

Il resta près de sa mère.

\scenebreak

Les jours suivants furent différents de tous ceux que Kamatsu avait connus.

Il n'y avait plus de route qui défilait sous ses pieds, plus de chariots à
visiter ni de nouveau paysage derrière le prochain virage. Ses parents et lui
restaient enfermés avec d'autres prisonniers, dans un espace dont Kamatsu
apprit rapidement chaque limite.

Il posa beaucoup de questions au début. Quand allaient-ils repartir ? Pourquoi
les Orcs les gardaient-ils ? Qu'est-ce qu'ils voulaient ? Est-ce que quelqu'un
viendrait les chercher ? Ses parents répondaient lorsqu'ils le pouvaient.

Souvent, ils ne le pouvaient pas.

Alors Kamatsu observa.

Il regarda les gens passer au-delà de leur prison. Il essaya de reconnaître les
Orcs qu'il avait vus pendant l'attaque et écouta leurs conversations sans en
comprendre les mots. Lorsqu'il n'avait rien d'autre à faire, il suivait du
regard les ombres qui se déplaçaient sur le sol ou cherchait des formes dans
les morceaux de ciel qu'il pouvait apercevoir.

Ses parents, eux, observaient autre chose. Kamatsu remarquait leurs voix plus
basses lorsqu'ils pensaient qu'il n'écoutait pas. Il voyait son père examiner
les alentours et sa mère suivre des yeux chaque Orc qui s'approchait. Il
comprenait qu'ils avaient peur.

Cela suffisait pour qu'il sache qu'il devait avoir peur lui aussi.

Il ne savait simplement pas encore de quoi.

Un jour, ce fut un enfant qui leur apporta à manger.

Kamatsu le remarqua immédiatement. Il avait vu des enfants orcs dans le
village, mais celui-ci s'était approché des prisonniers avec leur nourriture.
Il était à peine plus âgé que lui, ou peut-être simplement plus grand. Kamatsu
eut du mal à le déterminer.

Sa mère posa aussitôt une main sur son épaule. L'enfant orc déposa ce qu'il
avait apporté tandis que Kamatsu l'observait. Lorsqu'il releva les yeux vers
son visage, l'Orc le regardait aussi.

Pendant quelques secondes, aucun des deux ne parla.

Puis l'enfant repartit. Kamatsu suivit sa silhouette jusqu'à ce qu'elle
disparaisse.

— Il a quel âge ?

Son père tourna la tête vers lui.

— Qui ?

— Lui.

Kamatsu montra l'endroit par lequel l'enfant venait de partir.

— Je ne sais pas.

— Tu crois qu'il habite ici ?

Son père eut un bref regard vers sa mère.

— Probablement.

Kamatsu réfléchit à cette réponse. Jusqu'ici, les Orcs avaient surtout été
pour lui ceux qui avaient attaqué la caravane, ceux qui les avaient conduits
jusqu'au village et ceux qui les gardaient enfermés.

Il n'avait jamais pensé à ceux qui y étaient simplement nés.

\scenebreak

L'enfant orc revint le lendemain avec la nourriture.

Kamatsu le reconnut immédiatement. Il l'observa distribuer les rations aux
prisonniers et attendit qu'il arrive à leur hauteur avant de s'approcher de la
clôture. Sa mère posa une main sur son épaule.

— Reste près de nous.

Kamatsu acquiesça, mais fit tout de même un pas vers l'Orc lorsqu'il déposa
leur nourriture. Il regarda de nouveau la longue dent qui traversait son
oreille.

— Ça vient d'un animal ?

L'Orc leva les yeux vers lui.

— Recule.

Kamatsu obéit. Il retourna auprès de ses parents avec sa ration, non sans jeter
un dernier regard vers la dent que l'Orc portait à l'oreille avant que celui-ci
ne s'éloigne.

Le lendemain, Kamatsu l'attendit. Lorsqu'il le vit arriver avec les rations, il
resta près de ses parents jusqu'à ce que l'Orc atteigne leur partie de
l'enclos, puis s'approcha de nouveau.

— Tu as quel âge ?

L'Orc posa un bol.

— Recule.

— Moi, j'en ai dix.

Cette fois, l'Orc tourna la tête vers lui. Kamatsu attendit une réponse qui ne
vint pas et finit par retourner auprès de ses parents.

Cela ne l'empêcha pas de recommencer le jour suivant.

Les questions changeaient à chaque visite. Kamatsu voulait savoir ce qu'était
la dent qu'il portait à l'oreille, ce qu'il faisait dans le village ou ce qui
se trouvait dans la forêt alentour. L'Orc, lui, répondait presque toujours de
la même manière.

— Recule.

Kamatsu reculait, puis revenait avec une autre question le lendemain.

Ses parents appréciaient beaucoup moins cette habitude. Sa mère le rappelait
régulièrement lorsqu'il s'approchait de la clôture et son père lui demanda plus
d'une fois de laisser l'Orc tranquille. Kamatsu obéissait, au moins jusqu'à la
distribution suivante. Ils ne lui interdirent pourtant jamais complètement de
lui parler.

Kamatsu voyait bien qu'ils surveillaient chacune de ces rencontres. Il avait
aussi remarqué la façon dont son père suivait du regard les autres Orcs lorsque
l'un d'eux approchait de l'enclos. Avec cet enfant, cependant, rien ne s'était
jamais produit. Il apportait la nourriture, ordonnait à Kamatsu de reculer et
repartait.

Alors ses parents restaient attentifs sans intervenir davantage.

Quelques jours passèrent ainsi.

Un matin, l'Orc arriva avec les rations comme à son habitude. Kamatsu
s'approcha de la clôture et attendit qu'il lui tende un bol entre deux rondins.

Cette fois, il ne repartit pas immédiatement.

— Comment t'appelles-tu ?

L'Orc resta silencieux. Kamatsu s'attendait déjà à entendre la réponse
habituelle.

Elle ne vint pas.

L'Orc le regarda encore quelques secondes.

— Khaer Ghal.

Kamatsu sourit.

— Moi, c'est Kamatsu.

Khaer Ghal hocha légèrement la tête avant de reprendre sa distribution.
Kamatsu retourna auprès de ses parents avec son bol.

Cette fois, il avait obtenu une réponse.

\scenebreak

Le lendemain, Khaer Ghal avait presque terminé sa distribution lorsque Kamatsu
s'approcha de la clôture.

— Khaer Ghal !

L'Orc s'arrêta et se retourna.

— Quoi ?

Kamatsu resta un instant silencieux. Il ne s'était pas vraiment attendu à ce
que cela fonctionne.

— Rien.

Khaer Ghal fronça les sourcils.

— Alors pourquoi m'appelles-tu ?

Kamatsu réfléchit très sérieusement.

— Pour voir si tu répondrais.

Khaer Ghal le fixa quelques secondes.

— C'est idiot.

— Mais tu as répondu.

L'Orc resta silencieux, puis repartit. Kamatsu le regarda s'éloigner avec un
sourire.

\scenebreak

À partir de ce jour, quelque chose changea lentement entre les deux garçons.

Khaer Ghal ne devint pas soudainement bavard. Certaines questions de Kamatsu
ne recevaient qu'un haussement d'épaules ou un regard perplexe, et il arrivait
encore à l'Orc de terminer sa distribution sans lui adresser un mot. Mais il ne
repartait plus systématiquement dès que sa tâche était terminée.

Kamatsu en profita. Il voulait savoir pourquoi les Orcs vivaient sur cette
colline ou ce qu'il y avait dans les bois. Il l'interrogea aussi sur la longue
dent qui traversait son oreille, question qui revint suffisamment souvent pour
que l'Orc finisse par céder.

— Tu l'as combattue tout seul ?

— Oui.

— À huit ans ?

— Oui.

Kamatsu regarda son oreille.

— C'est elle qui t'a fait ça ?

Khaer Ghal porta les doigts vers le morceau de cartilage manquant.

— Elle m'a mordu.

— Et après ?

— Je l'ai tuée.

Kamatsu observa la dent quelques secondes.

— Et tu l'as gardée ?

Cette fois, quelque chose qui ressemblait à de la fierté apparut sur le visage
de Khaer Ghal.

— C'était ma première épreuve.

Kamatsu réfléchit. Quelques jours plus tôt, l'immense fosse au centre du
village lui avait paru étrange et impressionnante. Il savait maintenant que
Khaer Ghal y était descendu alors qu'il avait son âge pour affronter seul une
créature capable de lui arracher une partie de l'oreille.

Il ne savait toujours pas quoi penser de cet endroit.

Quelques jours plus tard, alors que Khaer Ghal récupérait les récipients vides,
Kamatsu désigna les autres prisonniers.

— Pourquoi nous gardez-vous ici ?

Khaer Ghal haussa les épaules.

— Pour travailler.

— Et après ?

L'Orc fronça les sourcils.

— Après quoi ?

— Quand on aura travaillé.

Khaer Ghal resta quelques secondes sans répondre.

— Certains travailleront encore. D'autres iront dans l'arène.

Kamatsu cessa de jouer avec le bord du récipient qu'il tenait.

— Pour quoi faire ?

— Combattre.

— Contre quoi ?

— Ça dépend. Un guerrier. Une créature.

Kamatsu attendit la suite, mais Khaer Ghal semblait avoir terminé son
explication.

— Et si on ne veut pas ?

L'Orc fronça de nouveau les sourcils.

— Pourquoi ne voudrais-tu pas ?

Kamatsu le regarda, surpris par la question.

— Parce que je pourrais mourir.

Khaer Ghal désigna la dent qui traversait son oreille.

— Moi aussi, j'aurais pu mourir.

— Mais toi, tu voulais y aller ?

Khaer Ghal hésita.

— C'était mon épreuve.

— Ce n'est pas ce que j'ai demandé.

Cette fois, le silence dura plus longtemps.

Kamatsu attendit. Khaer Ghal regarda un instant l'arène, puis la dent qu'il
portait à l'oreille. Lorsqu'il reporta son attention sur lui, il n'avait
toujours pas répondu.

— Chez moi, on ne fait pas ça, reprit Kamatsu.

Khaer Ghal releva les yeux.

— Alors vous faites quoi ?

Kamatsu sourit. Cette question-là lui plaisait beaucoup plus.

Il lui raconta les routes. Il lui parla des villages traversés avec ses
parents, des marchés où les voix se mêlaient aux odeurs de nourriture et des
villes dans lesquelles il aurait été incapable de compter toutes les maisons.
Il lui parla des chariots, des voyageurs rencontrés en chemin et des histoires
qu'on racontait le soir autour des feux.

Khaer Ghal l'interrompait souvent.

— Combien de maisons ?

— Je ne sais pas.

— Cent ?

Kamatsu rit.

— Beaucoup plus.

— Deux cents ?

— Khaer Ghal, je ne les ai pas comptées.

Une autre fois, Kamatsu lui parla de la mer.

— C'est comme un lac ? demanda Khaer Ghal.

— Beaucoup plus grand.

— Jusqu'où ?

— Jusqu'à ce que tu ne voies plus rien.

Khaer Ghal le regarda avec méfiance.

— Ça n'existe pas.

— Si.

— Tu l'as vue ?

Kamatsu hésita.

— Non.

— Alors comment sais-tu que ça existe ?

— Mon père l'a vue.

Khaer Ghal considéra cette réponse quelques secondes.

\scenebreak

Avec le temps, leurs conversations cessèrent de se limiter aux distributions.

Au début, Kamatsu aperçut simplement Khaer Ghal repasser près de l'enclos en
fin de journée. L'Orc ne s'arrêtait que quelques instants avant de repartir.
Puis ces passages devinrent plus fréquents et, lorsque le village retrouvait
son calme, il lui arriva de venir s'asseoir de l'autre côté de la barrière.

Kamatsu s'installait alors en face de lui. Il avait toujours quelque chose à
raconter : une anecdote arrivée sur la route, un endroit qu'il avait visité,
une mésaventure de la caravane ou une histoire qu'un voyageur lui avait
autrefois confiée. Khaer Ghal posait beaucoup de questions. Certaines
obligeaient Kamatsu à demander confirmation à ses parents, et son père
intervenait parfois pour corriger un détail que son fils avait arrangé à sa
manière.

Khaer Ghal revenait alors avec de nouvelles questions.

De son côté, il commença parfois à apporter autre chose que les rations
destinées aux prisonniers. Un morceau de nourriture supplémentaire apparaissait
avec le reste, puis un autre quelques jours plus tard. Kamatsu comprit
rapidement qu'ils venaient du propre repas de Khaer Ghal.

Kamatsu racontait une histoire. Khaer Ghal partageait ce qu'il avait réussi à
mettre de côté. Puis ils continuaient à parler jusqu'à ce que quelqu'un les
rappelle.

Une histoire contre quelques morceaux de nourriture.

Cela leur convenait parfaitement.

Les parents de Kamatsu continuaient d'observer ces rencontres avec prudence.
Au début, sa mère le gardait près d'elle chaque fois que Khaer Ghal
s'approchait, et son père ne quittait presque jamais l'Orc des yeux. Ils
avaient vu ce dont les guerriers du village étaient capables et Kamatsu
comprenait suffisamment bien leurs regards pour savoir qu'ils n'avaient pas
oublié.

Pourtant, Khaer Ghal n'avait jamais levé la main sur lui.

Surtout, lorsqu'il venait s'asseoir de l'autre côté de la barrière, Kamatsu
pouvait pendant un moment parler d'autre chose que de l'enclos, des Orcs ou de
ce qu'on ferait d'eux le lendemain. Il pouvait raconter un marché trop bruyant,
une ville immense ou une histoire entendue autour d'un feu, exactement comme il
l'aurait fait quelques jours plus tôt dans la caravane.

Ses parents ne l'encouragèrent jamais à aller vers Khaer Ghal, mais ils
cessèrent de le retenir.

Ils laissèrent simplement Kamatsu continuer à être Kamatsu.

Peu à peu, lui aussi cessa de voir seulement en Khaer Ghal l'un des Orcs du
village. Il remarquait désormais d'autres choses.

Khaer Ghal posait de plus en plus de questions sur l'extérieur, mais répondait
de moins en moins vite à celles qui concernaient son propre village. Lorsque
Kamatsu lui demandait pourquoi les prisonniers devaient travailler ou pourquoi
quelqu'un devait combattre dans l'arène s'il ne le voulait pas, l'Orc avait
d'abord toujours une réponse.

Parce que c'était ainsi.

Parce que cela se faisait.

Parce que cela avait toujours été comme ça.

Puis les réponses devinrent moins immédiates. Parfois, Khaer Ghal restait
simplement silencieux.

Une dizaine de jours s'était écoulée depuis leur arrivée, peut-être davantage.
Kamatsu ne les comptait plus vraiment.

Il savait seulement que Khaer Ghal reviendrait.

Et qu'il aurait probablement encore des questions.